Multivariate time-series data generated by sensor networks for environmental monitoring, industrial control, and building management embed critical operational intelligence and pose distinct analytical challenges that drive recent advances in time-series analysis~\cite{kumar2022iaq_review, szczurek2020iaq}. Among these challenges, anomaly detection on multivariate time series is a central problem with applications spanning fault diagnosis, structural health monitoring, and process safety. Indoor air quality (IAQ) monitoring exemplifies the multivariate sensor-stream regime: CO\textsubscript{2}, volatile organic compounds (VOC), temperature, humidity, and particulate matter exhibit complex co-variation patterns while each sensor type suffers from distinct noise characteristics, calibration drift rates, and failure modes~\cite{coleman_sensing_2018, abid2021iaq_deep}. Time-series anomaly detection has been extensively surveyed~\cite{chandola2009survey, aggarwal2017outlier, pang2021deep_survey}, with recent IAQ-specific studies leveraging deep architectures for both forecasting and outlier identification~\cite{choi2024air_anomaly, wei2023lstm_ae_iaq}. In practice, deployment necessitates solutions that confront challenges unique to physical sensor networks, situated within a mathematically rigorous time-series framework—challenges that conventional time-series models fail to explicitly address.

This study is motivated by three sensor-specific challenges. First, multi-sensor deployments exhibit heterogeneous reliability across time. Individual sensors degrade at different rates: electrochemical VOC sensors drift over weeks, optical dust sensors saturate in high-concentration events, and temperature sensors maintain stable calibration for months~\cite{sanchez-ferrera_review_2025}. Standard attention mechanisms treat all time steps equally, yet predictions for a degrading VOC sensor at time $t$ are inherently less reliable than predictions for a stable temperature sensor at the same instant. A sensor-aware model should dynamically down-weight attention to unreliable channels. Second, anomaly scoring must adapt to context. Most approaches combine reconstruction error and prediction uncertainty using fixed weights~\cite{ghimire_electricity_2025, cui_sde-hnn_2025}, but the relative informativeness of these signals depends on sensor state: during stable operation, large reconstruction errors reliably indicate anomalies; During degradation events, a high level of uncertainty may actually serve as a more reliable signal, even though reconstruction errors appear moderate. A fixed-weight scoring system combines these two distinct conditions, thereby reducing detection accuracy during sensor transitions. Third, anomaly detectors should produce actionable output. Binary flagging without sensor-level attribution leaves facility operators unable to determine whether a detected anomaly reflects a genuine environmental hazard, a degrading CO\textsubscript{2} sensor requiring recalibration, or a VOC spike from cleaning activities~\cite{xu2022anomalytransformer}. For anomaly detection to inform maintenance decisions, it must identify which sensors contribute to each anomaly and provide uncertainty-aware estimates rather than only point predictions.

The Temporal Convolutional Network (TCN) architecture~\cite{bai2018tcn} provides an effective backbone for sequence modeling through dilated causal convolutions, matching or exceeding recurrent models on time-series benchmarks. Monte Carlo (MC) dropout~\cite{gal2016dropout_bayesian} enables uncertainty estimation from deterministic architectures without requiring explicit Bayesian inference. Neither technique alone, however, addresses the sensor-specific challenges described above. The principal contribution of this study lies in the integration of uncertainty signals, not as independent outputs, but as internal mechanisms that actively influence attention computation, anomaly scoring, and sensor-level attribution, thereby aligning with the operational realities of physical sensor deployments.

We present an uncertainty-aware TCN framework for multivariate sensor anomaly detection, making three contributions tightly coupled to the challenges above:

\begin{enumerate}[leftmargin=*, nosep]
    \item \textbf{Adaptive Uncertainty-Aware Attention (AUAA).} An attention mechanism that modulates temporal attention weights using per-sensor predictive uncertainty estimates via an uncertainty gate $g_t = 1/(1 + \sigma_t)$. This makes attention computation sensitive to sensor-specific reliability fluctuations, enabling the model to down-weight unreliable channels.
    
    \item \textbf{Dynamic Weight Adapter for Hybrid Anomaly Scoring.} A learnable gating module that contextually balances reconstruction error against predictive uncertainty using a GRU over recent weight history. This replaces static combination weights with a mechanism that adapts to sensor operating conditions.
    
    \item \textbf{Feature-Level Anomaly Attribution with Theoretical Bounds.} A feature-importance learner providing per-sensor explanations, complemented by Chebyshev-based uncertainty bounds that guarantee $P(\text{false positive}) \leq 1/k^2$. This delivers actionable, sensor-specific explanations with provable calibration guarantees.
\end{enumerate}

% We validate through four-month IAQ sensor data, comprehensive ablation studies isolating each component, cross-dataset evaluation on the Skoltech Anomaly Benchmark (SKAB)~\cite{skab_dataset}.
% a publicly available labelled industrial water-circulation corpus whose ground-truth labels are exogenous to our model development pipeline and sensitivity analysis of key hyperparameters.